\section{Related Work}
\label{sec:related}

\subsection{AI-Assisted Academic Writing}

Recent advances in large language models~\cite{openai2023gpt4,anthropic2024claude}
have spurred interest in AI-assisted scholarly writing. General-purpose chat
interfaces such as ChatGPT, Claude, and Gemini provide text generation
capabilities but lack project-level awareness of manuscript structure, citation
databases, and compilation requirements.

Specialized tools address specific aspects of the writing process.
Writefull~\cite{writefull} provides language polishing and grammar checking
for academic text through a dedicated interface. Paperpal~\cite{paperpal} offers
comprehensive manuscript preparation assistance, including language checks and
reference formatting. Trinka~\cite{trinka} specializes in technical and academic
language correction. However, these tools operate on individual passages rather
than managing complete manuscript projects with source files, bibliographies,
and compilation.

Research on LLM-assisted writing has explored prompt engineering for academic
genres, fine-tuning for scientific domains, and evaluation of generated text
quality. Liang et al.~\cite{liang2024llmwriting} conducted a systematic survey
of LLM use across the academic writing pipeline, identifying tasks ranging from
ideation to revision. They found that while LLMs are widely used for drafting
and editing, systematic integration with document engineering infrastructure
remains limited. These studies evaluate text quality
in isolation rather than the end-to-end workflow of producing compilable,
verifiable scholarly documents---a gap that \software{} aims to fill.
A further challenge is that preprint repositories such as arXiv assign DOIs
(\texttt{10.48550/arXiv.\ldots}) that CrossRef does not index, requiring
dedicated API integration that current writing tools do not provide.

\subsection{Collaborative \LaTeX{} Editing}

Online \LaTeX{} editing platforms dominate collaborative scholarly writing.
Overleaf~\cite{overleaf} provides real-time collaboration, version history,
and integrated compilation, serving over 10 million users globally.
Overleaf has recently introduced AI-assisted features~\cite{overleaf2024ai},
including grammar suggestions and summarization. Yet its cloud-centric
architecture means manuscripts reside on remote servers, and AI integration is
constrained by the platform's service model. Researchers concerned about data
privacy or requiring offline access must seek alternatives.

Local \LaTeX{} editors such as TeXstudio~\cite{texstudio}, VS Code with the
\LaTeX{} Workshop extension, and Emacs with AUCTeX provide powerful editing
environments with syntax highlighting, code completion, and local compilation.
While cloud platforms offer real-time collaboration out of the box, local
setups provide superior compilation performance, offline availability, and
full control over \TeX{} distributions---at the cost of manual environment
configuration. These tools increasingly support AI integration through
extensions, but none provide a unified, project-aware AI workspace with built-in
verification and workflow orchestration.

Modern reproducible-authoring tools address a related but distinct problem.
R Markdown~\cite{rmarkdown} and Quarto~\cite{quarto} combine prose, executable
code cells, citations, and multi-format rendering, making them strong choices
for computational notebooks, statistical reports, and literate data analysis.
Their primary abstraction is the executable document: code and narrative are
woven into HTML, PDF, Word, or presentation outputs. Quarto, the successor to
R Markdown, provides a unified authoring format that targets multiple output
engines---\LaTeX{}, Typst, and revealjs---and integrates with version control,
project-level metadata (\texttt{\_quarto.yml}), and CI/CD pipelines. RStudio
(also from Posit, the Quarto developer) provides an IDE that tightly couples
code execution, notebook preview, and package management for R and Python
workflows.

\software{} differs from Quarto and RStudio Markdown in three fundamental ways.
First, the primary artifact: Quarto and R Markdown treat the \texttt{.qmd}/\texttt{.Rmd}
source file as the unit of authoring, while \software{} treats the manuscript
directory itself as the primary artifact, preserving ordinary \LaTeX{} project
structure with separate \texttt{.tex}, \texttt{.bib}, and \texttt{.sty} files.
Second, the role of AI: Quarto and RStudio do not provide built-in AI
assistance with permission-aware interaction modes; AI integration in these
ecosystems typically relies on external extensions (e.g., GitHub Copilot in
RStudio) that operate at the text-buffer level without project-level awareness
of citation databases or compilation state. \software{} places AI assistance
behind explicit permission modes (Chat, Agent, Tools) and integrates it with
project-level services such as citation verification and compilation. Third,
the verification model: Quarto focuses on reproducible rendering from code
chunks, while \software{} focuses on verifying that AI-generated edits preserve
compilability, citation accuracy, and document integrity. Thus, Quarto and
R Markdown emphasize reproducible publishing from code-rich documents, whereas
\software{} emphasizes controlled AI assistance, citation verification, and
project-level document engineering for existing \LaTeX{} manuscript workflows.
These approaches are complementary rather than competing: a researcher could
use Quarto for computational reproducibility and \software{} for AI-assisted
manuscript engineering within the same project.

\subsection{Agent-Based Writing Systems}

The concept of AI agents---systems that can plan, use tools, and execute
multi-step tasks---has been applied to writing. Microsoft's Research
Copilot~\cite{microsoft2024researchcopilot} and related tools augment
research workflows with AI capabilities, including literature review,
data analysis, and writing assistance. LangChain~\cite{langchain},
AutoGen~\cite{autogen}, CrewAI~\cite{crewai}, and
LangGraph~\cite{langgraph} enable tool-using or multi-agent workflows where
different agents handle planning, drafting, reviewing, and execution.

However, these systems typically operate on plain text and lack direct
integration with \LaTeX{} project structures, citation databases, and
compilation pipelines. Their abstractions are general-purpose chains, agents,
tools, and conversations; the manuscript-specific responsibilities---preserving
cross-references, checking BibTeX keys, compiling \LaTeX{}, and deciding whether
an edit should enter the source tree---must be assembled by the application
developer.

LangChain provides composable chains and agent executors that can orchestrate
multi-step LLM interactions with tool use. In an academic writing scenario, one
could build a LangChain agent that drafts sections, verifies citations via
custom tools, and compiles \LaTeX{} documents. However, LangChain's default
agent loop applies tool outputs automatically and continues execution without
mandatory human approval gates. The researcher becomes a post-hoc reviewer of
fully generated output rather than an active decision-maker at each step.
AutoGen extends this model by enabling multi-agent conversations where
different agents (e.g., a ``writer,'' a ``reviewer,'' an ``editor'') collaborate
autonomously. While this architecture can produce sophisticated outputs through
agent deliberation, it further reduces the human role: the researcher observes
agent-to-agent dialogue and intervenes only when the autonomous loop completes
or stalls. CrewAI adopts a role-based orchestration pattern where agents are
assigned specialized roles and goals, executing tasks in sequence or parallel.
All three frameworks share a common design philosophy: AI as an autonomous
executor that minimizes human intervention, with the human serving primarily as
a supervisor or approver of final results.

\software{} takes the opposite approach: \emph{AI as a proposer rather than a
replacement author.} Chat mode cannot modify files, Agent mode returns
reviewable diffs that require explicit accept/reject actions, and Tools mode is
explicitly scoped to the project directory with full audit logging. This
``proposer-not-replacer'' philosophy differs fundamentally from the
``autonomous executor'' paradigm of LangChain, AutoGen, and CrewAI in three
ways: (1) every file modification requires a deliberate human action; (2) the
system never auto-applies AI-generated changes and never continues execution
past a human checkpoint without explicit approval; and (3) the audit trail
records not only what the AI proposed but also whether the human accepted or
rejected each change. This design keeps scholarly authorship and final editorial
authority with the researcher while still allowing agentic tooling to assist
with bounded document engineering tasks. The tradeoff is reduced automation:
tasks that LangChain or AutoGen could complete autonomously (e.g., ``revise all
sections for clarity and recompile'') require human interaction at each
checkpoint in \software{}. We argue that this tradeoff is appropriate for
scholarly writing, where authorial accountability cannot be delegated to an
autonomous agent.

\subsection{Human-in-the-Loop AI Systems}

Human-in-the-loop (HITL) approaches have been extensively studied in machine
learning~\cite{hitl_survey}, where human feedback guides model training and
decision-making. In writing tools, HITL principles suggest that AI should
propose rather than impose changes, and that humans should retain decision
authority at key junctures.

Grammarly~\cite{grammarly} exemplifies HITL in writing tools: it highlights
issues and suggests changes, but the user applies them manually. GitHub
Copilot~\cite{copilot} applies similar principles for code generation---it shows
completions in-line, and the developer decides whether to accept. \software{}
generalizes this principle across the entire academic writing workflow: the
Chat mode enables free-form discussion without risk of changes, the Agent mode
shows proposed edits as inline diffs for review, and the Tools mode provides
controlled execution for multi-step operations within a sandboxed directory.

\subsection{Software Tools for Research Software Engineering}

Software: Practice and Experience has a tradition of publishing papers on
research software tools and platforms that emphasize practical engineering
contributions. Notable examples include
Jupyter~\cite{jupyter2015software} for interactive computing, the Git
ecosystem for version control~\cite{git}, and continuous integration
platforms for research software~\cite{githubactions}. These contributions share a
common emphasis on practical software engineering: tools designed for real
workflows, validated through use, and documented with attention to
implementation tradeoffs.

\software{} contributes to this tradition by providing a practical software
platform for AI-assisted scholarly writing, with emphasis on real-world design
decisions, implementation experience, and lessons learned from building a system
that bridges LLMs and document engineering. The paper follows the SPE
convention of reporting on an implemented system with attention to architecture,
testing methodology, performance characteristics, and deployment
considerations---rather than positioning the contribution primarily as a
theoretical advance or a user study.
